% =============================================================================
\chapter{Tło i prace pokrewne}
\label{ch:background}
% =============================================================================

\section{Paradygmaty programowania}
\label{sec:paradigms}

Paradygmat programowania to sposób strukturyzowania obliczeń, który język wspiera i do którego
zachęca~\cite{floyd1979paradigms,vanroy2004ctm}. Rozważamy cztery klasyczne paradygmaty i tam,
gdzie to możliwe, wybieramy po dwa języki na paradygmat. Jeden z nich jest szeroko stosowany,
a drugi jest mniej popularnym, „czystszym” przedstawicielem. Pozwala to oddzielić wpływ paradygmatu
od wpływu popularności pojedynczego języka.

\paragraph{Imperatywny (proceduralny): C i Pascal} Programy są ciągami instrukcji modyfikujących
stan. Dane to rekordy i tablice, a przepływ sterowania opiera się na pętlach i instrukcjach
warunkowych. W C pamięcią zarządza się ręcznie za pomocą \code{malloc}/\code{free}, więc wycieki
i niepoprawne dostępy stanowią odrębną klasę defektów. Pascal wprowadza silniejsze typowanie
i zarządzane tablice dynamiczne, ale zachowuje \code{New}/\code{Dispose} dla wskaźników.

\paragraph{Obiektowy: Java i Smalltalk} Zachowanie znajduje się w obiektach, które hermetyzują stan
i komunikują się za pomocą komunikatów. Java jest typowana statycznie i ma bogatą bibliotekę
standardową. Smalltalk to typowany dynamicznie język, w którym „wszystko jest obiektem”, a nawet
struktury sterujące (\code{ifTrue:}, \code{do:}) są komunikatami wysyłanymi do obiektów i bloków.

\paragraph{Funkcyjny: Haskell i Scala} Obliczenie polega na wartościowaniu wyrażeń bez efektów
ubocznych. Haskell jest czysty i leniwy. Brak wartości i błędy modeluje algebraicznymi typami danych
(\code{Maybe}, \code{Either}) i preferuje zwinięcia, dopasowanie wzorców oraz trwałe struktury
danych. Scala łączy cechy funkcyjne i obiektowe na JVM, oferując niemodyfikowalne kolekcje, klasy
przypadków (ang.~\emph{case classes}) i \code{Option}.

\paragraph{Logiczny: Prolog} Program jest zbiorem relacji (faktów i reguł), a obliczenie polega na
poszukiwaniu dowodu za pomocą unifikacji i nawracania. Idiomatyczny Prolog używa głów klauzul do
analizy przypadków, akumulatorów do iteracji oraz porażki (a nie wyjątków) do wyrażenia „braku
odpowiedzi”. Odcięć i stanu globalnego używa się oszczędnie.

\section{Benchmarki generowania kodu}
\label{sec:related-benchmarks}

\paragraph{Benchmarki na poziomie funkcji} HumanEval~\cite{chen2021codex} zawiera 164 ręcznie
napisane problemy w Pythonie, z których każdy składa się z sygnatury funkcji, docstringu i średnio
7,7 testu jednostkowego. Wprowadził metrykę pass@$k$. MBPP~\cite{austin2021mbpp} zawiera 974 krótkie
problemy w Pythonie zebrane metodą crowdsourcingu. Oba stały się domyślnymi miarami dla modeli
kodu. \citet{liu2023evalplus} pokazali, że ich testy są często zbyt słabe: rozszerzenie HumanEval
o wiele dodatkowych testów (HumanEval+) wyraźnie obniża raportowane wyniki, ponieważ rozwiązania
wyglądające wiarygodnie, lecz błędne, przestają być akceptowane.

\paragraph{Benchmarki wielojęzyczne} MultiPL-E~\cite{cassano2023multiple} tłumaczy problemy i testy
z HumanEval i MBPP na 18 kolejnych języków za pomocą kompilatorów z Pythona do każdego z języków
docelowych. HumanEval-X~\cite{zheng2023codegeex} i MBXP~\cite{athiwaratkun2023mbxp} realizują tę
samą ideę. Benchmarki te pokazały, że wyniki spadają wraz z ilością danych treningowych dostępnych
dla danego języka~\cite{cassano2023multiple,cassano2024knowledge}. Problemy pozostają jednak
przetłumaczonymi problemami z Pythona. Są to krótkie, bezstanowe funkcje na listach, napisach
i liczbach, które rzadko wymagają projektowania specyficznego dla paradygmatu. Programowanie
logiczne jest w nich niemal nieobecne.

\paragraph{Większe jednostki i kontaminacja} ClassEval~\cite{du2023classeval} ocenia generowanie
całych klas, a SWE-bench~\cite{jimenez2024swebench} rozwiązywanie rzeczywistych zgłoszeń z GitHuba.
LiveCodeBench~\cite{jain2024livecodebench} z czasem dobiera nowe problemy z konkursów, aby uniknąć
kontaminacji, czyli obecności problemów z benchmarku i ich rozwiązań w danych treningowych, co
zawyża wyniki popularnych benchmarków~\cite{riddell2024contamination}.

Nasz benchmark sytuuje się między tymi nurtami. Jego zadania są większe niż problemy HumanEval
(kontrakty z wieloma funkcjami, struktury danych, własność pamięci), ale wystarczająco małe, by
zaimplementować je od zera w siedmiu językach z \emph{identycznymi} testami. Według naszej wiedzy
żaden istniejący benchmark nie utrzymuje jednocześnie stałego zadania i stałych testów we wszystkich
czterech klasycznych paradygmatach ani nie ocenia przy tym idiomatyczności wyniku.

\section{Samonaprawa z informacją zwrotną z wykonania}
\label{sec:related-repair}

Gdy wygenerowany kod zawodzi, można pokazać modelowi błąd i poprosić o poprawkę.
Self-Debugging~\cite{chen2024selfdebug} uczy modele wyjaśniać i naprawiać swój kod na podstawie
wyników wykonania, a Reflexion~\cite{shinn2023reflexion} pozwala agentowi przechowywać werbalną
informację zwrotną między próbami. \citet{olausson2024selfrepair} stwierdzili, że korzyść
z samonaprawy jest rzeczywista, lecz mniejsza, niż się często zakłada. Ogranicza ją przede wszystkim
jakość informacji zwrotnej, a tę kontroluje właśnie program uruchamiający testy. Nasz system
realizuje samonaprawę z informacją zwrotną z wykonania w ramach jednej konwersacji. Model widzi
komunikaty kompilatora, logi testów i nazwy niezaliczonych testów, ale nigdy kodu testów
(\cref{sec:repair-loop}).

\section{LLM jako sędziowie}
\label{sec:related-judge}

Użycie silnego modelu LLM do oceniania wyników (ang.~\emph{LLM-as-a-judge}) skaluje ewaluację tam,
gdzie nie istnieje odpowiedź wzorcowa. \citet{zheng2023judging} raportują zgodność między silnymi
sędziami LLM a ludźmi porównywalną ze zgodnością między ludźmi. Dokumentują też obciążenia związane
z pozycją, rozwlekłością i preferowaniem własnych odpowiedzi. \citet{panickssery2024selfpreference}
pokazują, że ewaluatory LLM potrafią rozpoznawać i faworyzować własne generacje. W przypadku kodu
ICE-Score~\cite{zhuo2024icescore} używa LLM do oceny użyteczności i poprawności funkcjonalnej bez
wykonywania kodu. My używamy sędziego wyłącznie do tego, czego testy nie mogą zmierzyć, czyli
idiomatyczności, czytelności i projektu, i ograniczamy go rubryką właściwą dla języka,
ustrukturyzowaną odpowiedzią w formacie JSON i zaślepieniem (\cref{sec:judge}).

\section{Języki programowania a jakość kodu}
\label{sec:related-quality}

Spór o to, czy wybór języka wpływa na jakość oprogramowania, toczył się na długo przed pojawieniem
się LLM. \citet{ray2014largescale} raportowali niewielkie, lecz istotne związki między językami
a commitami naprawiającymi defekty na GitHubie. Późniejsza replikacja \citet{berger2019impact}
wykazała, że większość tych efektów znika po skorygowaniu danych i analizy statystycznej. Ta historia
skłania do ostrożności: efekty języka łatwo pomylić z efektami projektu, danych i pomiaru, a ich
badanie wymaga starannej statystyki. Wracamy do tej kwestii w \cref{sec:threats}.
